% Part: incompleteness
% Chapter: incompleteness-provability
% Section: second-incompleteness-thm

\documentclass[../../../include/open-logic-section]{subfiles}

\begin{document}

\olfileid{inc}{inp}{2in}

\olsection{مبرهنة عدم الاكتمال الثانية}

كيف يمكننا التعبير عن القول إن~$\Th{PA}$ لا تثبت اتساقها؟ إن القول
إن~$\Th{PA}$ غير متسقة يكافئ القول إن
$\Th{PA} \Proves \eq[0][1]$. ولذلك يمكننا أن نتخذ عبارة الاتساق
$\OCon[\Th{PA}]$ هي الـ!!{sentence}~$\lnot
\OProv[\Th{PA}](\gn{\eq[0][1]})$، وعندئذ تفي المبرهنة الآتية بالغرض:

\begin{thm}
\ollabel{thm:second-incompleteness}
إذا افترضنا أن~$\Th{PA}$ متسقة، فإن~$\Th{PA}$ لا !!{derive}
$\OCon[\Th{PA}]$.
\end{thm}

من المهم ملاحظة أن المبرهنة تعتمد على التمثيل المعين لـ
$\OCon[\Th{PA}]$ (أي التمثيل المعين لـ$\OProv[\Th{PA}](y)$). وكل ما
سنستعمله هو أن تمثيل~$\OProv[\Th{PA}](y)$ يحقق شروط
!!{derivability} الثلاثة، ولذلك تتعمم المبرهنة على كل نظرية لها
!!a{derivability} محمول يحقق هذه الخواص.

ومن المفيد قراءة مخطط غودل للحجة، إذ تأتي المبرهنة كخاتمة موفقة. تسير
الحجة على النحو الآتي. لتكن~$!G_\Th{PA}$ جملة غودل التي بنيناها في
برهان~\olref[1in]{thm:first-incompleteness}. لقد بيّنا: «إذا كانت
$\Th{PA}$ متسقة، فإن~$\Th{PA}$ لا !!{derive} $!G_\Th{PA}$». وإذا
صغنا هذا صوريًا \emph{داخل}~$\Th{PA}$، كان لدينا برهان على
\[
\OCon[\Th{PA}] \lif \lnot \OProv[\Th{PA}](\gn{!G_\Th{PA}}).
\]
افترض الآن أن~$\Th{PA}$ !!{derive}s $\OCon[\Th{PA}]$. عندئذ
!!{derive}s أيضًا $\lnot \OProv[\Th{PA}](\gn{!G_\Th{PA}})$. ولكن لما
كانت~$!G_\Th{PA}$ جملة غودل، كان هذا مكافئًا لـ$!G_\Th{PA}$. ومن ثم
!!{derive}s $\Th{PA}$ الصيغة~$!G_\Th{PA}$.

لكننا نعلم أنه إذا كانت~$\Th{PA}$ متسقة، فهي لا !!{derive}
$!G_\Th{PA}$!{} ولذلك، إذا كانت~$\Th{PA}$ متسقة، فلا يمكنها أن
!!{derive} $\OCon[\Th{PA}]$.

ولزيادة الحجة ضبطًا، سندع~$!G_\Th{PA}$ تكون جملة غودل لـ$\Th{PA}$،
ونستعمل شروط~!!{derivability} (P1)--(P3) لنبين أن~$\Th{PA}$
!!{derive}s $\OCon[\Th{PA}] \lif !G_\Th{PA}$. وسيبين هذا أن
$\Th{PA}$ لا !!{derive} $\OCon[\Th{PA}]$. وفيما يأتي مخطط البرهان
داخل~$\Th{PA}$. (وللتبسيط نحذف المؤشرات السفلية~$\Th{PA}$.)
\begin{align}
& !G \liff \lnot \OProv(\gn{!G}) \ollabel{G2-1}\\
& \qquad\text{إن~$!G$ جملة غودل}\notag \\
& !G \lif \lnot \OProv(\gn{!G}) \ollabel{G2-2}\\
  & \qquad\text{من \olref{G2-1}} \notag\\
& !G \lif
  (\OProv(\gn{!G}) \lif \lfalse) \ollabel{G2-3}\\
  & \qquad\text{من \olref{G2-2} بالمنطق}\notag\\
& \OProv(\gn{
    !G \lif
    (\OProv(\gn{!G}) \lif \lfalse)
  }) \ollabel{G2-4}\\
  & \qquad\text{من \olref{G2-3} بحسب الشرط P1} \notag\\
& \OProv(\gn{!G}) \lif
  \OProv(\gn{
    (\OProv(\gn{!G}) \lif \lfalse)
    }) \ollabel{G2-5}\\
  & \qquad\text{من \olref{G2-4} بحسب الشرط P2} \notag\\
& \OProv(\gn{!G}) \lif (\OProv(\gn{\OProv(\gn{!G})}) \lif \OProv(\gn{\lfalse})) \ollabel{G2-6}\\
  & \qquad\text{من \olref{G2-5} بحسب الشرط P2 والمنطق} \notag\\
& \OProv(\gn{!G}) \lif
  \OProv(\gn{\OProv(\gn{!G})}) \ollabel{G2-7}\\
   & \qquad\text{بحسب P3} \notag\\
& \OProv(\gn{!G}) \lif \OProv(\gn{\lfalse}) \ollabel{G2-8}\\
  & \qquad \text{من \olref{G2-6} و\olref{G2-7} بالمنطق}\notag\\
& \OCon \lif \lnot \OProv(\gn{!G}) \ollabel{G2-9}\\
  & \qquad\text{بعكس نقيض \olref{G2-8} و$\OCon \ident \lnot \OProv(\gn{\lfalse})$}\notag \\
& \OCon \lif !G \notag\\
  & \qquad\text{من \olref{G2-1} و\olref{G2-9} بالمنطق}\notag
\end{align}
ولا ينطوي استعمال المنطق أعلاه إلا على حقائق أولية من المنطق
القضوي؛ فمثلًا تستعمل~\olref{G2-3} العلاقة
$\Proves \lnot!A \liff (!A\lif \lfalse)$، وتستعمل~\olref{G2-8}
$!A \lif (!B \lif !C), !A \lif !B \Proves !A \lif !C$. أما استعمال
الشرط~P2 في~\olref{G2-5} و\olref{G2-6} فيعتمد على حالتين من~P2، هما
$\OProv(\gn{!A \lif !B}) \lif
(\OProv(\gn{!A}) \lif \OProv(\gn{!B}))$. في الأولى
$!A \ident !G$ و$!B \ident \OProv(\gn{!G}) \lif \lfalse$؛ وفي
الثانية $!A \ident \OProv(\gn{!G})$ و$!B \ident \lfalse$.

والصيغة الأكثر تجريدًا لمبرهنة عدم الاكتمال الثانية هي الآتية:

\begin{thm}
\ollabel{thm:second-incompleteness-gen} لتكن~$\Th{T}$ أي نظرية
متسقة و!!{axiomatized} تمتدّ~$\Th{Q}$، ولتكن
$\OProv[\Th{T}](y)$ أي صيغة تحقق شروط~!!{derivability} P1--P3
لـ$\Th{T}$. عندئذ لا !!{derive} $\Th{T}$ الصيغة~$\OCon[\Th{T}]$.
\end{thm}

\begin{prob}
بيّن أن~$\Th{PA}$ !!{derive}s $!G_{\Th{PA}} \lif \OCon[\Th{PA}]$.
\end{prob}

\begin{digress}
مغزى القصة أنه لا يمكن لأي نظرية متسقة «معقولة» للرياضيات أن
!!{derive} عبارة اتساقها. افترض أن~$\Th{T}$ نظرية في الرياضيات تتضمن
$\Th{Q}$ واستدلال هيلبرت «التناهي» (أيًا يكن معناه). عندئذ لا يمكن
لـ$\Th{T}$ كلها أن !!{derive} عبارة اتساق~$\Th{T}$، ومن باب أولى لا
يمكن لجزئها التناهي أن !!{derive} عبارة اتساق~$\Th{T}$ أيضًا. وبهذا
المعنى لا يمكن أن يوجد برهان تناهٍ على اتساق «الرياضيات كلها».

وثمة قدر من المرونة في تفسير مصطلح «تناهٍ». وقد أقر غودل في بحثه عام
1931 باحتمال أن يقع شيء نعدّه «تناهيًا» خارج أنواع الرياضيات التي أراد
هيلبرت صورنتها. غير أن غودل كان متسامحًا؛ فمن العسير اليوم أن نتصور
كيف نجد شيئًا يمكن تسميته تناهيًا على نحو معقول ولا يمكن صورنته في
$\Th{ZFC}$ مثلًا، أي نظرية مجموعات زرميلو--فرينكل مع بديهية الاختيار.
\end{digress}

\end{document}
